\chapter*{Introduzione}
\addcontentsline{toc}{chapter}{Introduzione}

\phantomsection
\section*{Obiettivi della tesi}
\addcontentsline{toc}{section}{Obiettivi della tesi}
La seguente tesi si propone di esplorare le architetture web moderne e le tecniche di ottimizzazione correlate, con
l'obiettivo di fornire una panoramica completa delle sfide e delle opportunità legate allo sviluppo di applicazioni web
efficienti e scalabili. In particolare, l'obiettivo principale è quello di analizzare le differenze tra le architetture
di rendering e le loro peculiarità, per fornire gli strumenti necessari per scegliere l'approccio più adatto in base
alle esigenze specifiche di un progetto.
Verranno inoltre esaminate le tecniche di hydration, code splitting, lazy loading e speculative loading, con l'obiettivo
di fornire una comprensione approfondita delle strategie di ottimizzazione delle performance e del caricamento delle
risorse in applicazioni web moderne.
Ma questo solo dopo aver introdotto le basi dello sviluppo web, con particolare attenzione al runtime environment di
JavaScript, alle Web APIs e al Document Object Model (DOM), per fornire un contesto completo per la comprensione delle
architetture e delle tecniche di ottimizzazione trattate.

\phantomsection
\section*{Importanza delle architetture web moderne}
\addcontentsline{toc}{section}{Importanza delle architetture web moderne}

Negli ultimi anni, lo sviluppo di applicazioni web ha subito un'evoluzione radicale, guidata dalla crescente necessità
di prestazioni elevate, interfacce utente reattive e una migliore esperienza utente. Le tradizionali pagine web
statiche, generate interamente dal server, non sono più sufficienti per soddisfare le esigenze moderne di
applicazioni dinamiche e interattive, come i social network, le piattaforme di e-commerce e i
servizi di streaming. Questo ha portato all’emergere di nuove architetture web che cercano di ottimizzare il modo in cui
i contenuti vengono generati e serviti, bilanciando le esigenze di performance, usabilità, e riduzione dei costi e della
difficoltà di sviluppo e manutenzione, in modi diversi a seconda delle esigenze del progetto.

Un'altra esigenza critica delle moderne architetture web è l'ottimizzazione del caricamento delle risorse, che ha
portato all'adozione di tecniche come il lazy loading, il code splitting e lo speculative loading. Questi approcci consentono di caricare solo
le parti dell'applicazione necessarie, riducendo il tempo di caricamento iniziale e migliorando l'esperienza utente,
soprattutto su dispositivi con risorse limitate.

In questo contesto, l'adozione di architetture flessibili e performanti è diventata essenziale per lo sviluppo di
applicazioni web che possano soddisfare gli standard moderni in termini di reattività, scalabilità e usabilità. Questo
elaborato intende esplorare tali architetture e le tecniche di ottimizzazione correlate, con l'obiettivo di fornire una
comprensione approfondita delle sfide e delle opportunità legate allo sviluppo di applicazioni web efficienti e
scalabili.

\phantomsection
\section*{Struttura della tesi}
\addcontentsline{toc}{section}{Struttura della tesi}
La tesi è strutturata in quattro capitoli principali, ciascuno dei quali affronta un aspetto specifico delle
architetture web moderne e delle tecniche di ottimizzazione correlate.

Il primo capitolo presenta le fondamenta dello sviluppo web, iniziando con una panoramica del JavaScript runtime
environment, analizzando il funzionamento del JavaScript engine, entrando nel dettaglio di come il codice viene eseguito
e ottimizzato, e spiegando il concetto di Execution Context e l'Event Loop. 
Vengono esaminate le Web APIs, un insieme di interfacce e metodi che permettono ai programmatori di interagire con il
browser e manipolare il DOM. Viene trattato il Document Object Model (DOM), una rappresentazione gerarchica dei
documenti HTML e XML, che consente di manipolare la struttura e il contenuto di una pagina web.
Vengono esaminate le caratteristiche fondamentali del runtime environment di JavaScript, con particolare attenzione alla gestione delle operazioni asincrone e alla gestione degli eventi.

Il secondo capitolo introduce il contesto generale delle architetture web moderne, analizzando le principali differenze tra Server-Side Rendering (SSR) e Client-Side Rendering (CSR), e presentando le architetture ibride come soluzione
intermedia. Vengono esaminate le caratteristiche, i vantaggi e gli svantaggi di ciascuna architettura, e vengono confrontate in termini di prestazioni, usabilità e indicizzazione SEO.

Il terzo capitolo approfondisce il concetto di hydration, una tecnica utilizzata per convertire pagine HTML statiche in
pagine web dinamiche post-rendering. Viene spiegato il funzionamento di questa tecnica, le due fasi principali dell processo di hydration, e le diverse implementazioni possibili di questa tecnica in applicazioni web moderne.

Il quarto capitolo si concentra sulle tecniche di ottimizzazione delle performance: il code splitting, il lazy
loading e lo speculative loading, che consentono di migliorare il caricamento delle risorse e l'esperienza utente.
Vengono esaminate le caratteristiche di queste tecniche, i vantaggi che offrono e le modalità di implementazione in
applicazioni web moderne.